% ==== 题名页、摘要、译者说明、目录 ====
\begin{abstract}
\noindent
在这篇教学性综述中，我们介绍并研究一种被称为陈--西蒙斯理论的三维拓扑场论。我们从 $\mathrm{U}(1))$ 作用量的现象学出发：讨论它的经典运动方程、它的量子化以及它的可观测对象。随后我们转向非阿贝尔理论，分别用路径积分与正则两种形式处理它。我们尤其关注流形与纽结的标架化，以及量子理论的几何结构。
\end{abstract}

% [Tutor Note: 物理图像] 原文摘要只有四句，却是全文骨架：U(1) 现象学 -> 量子化 -> 可观测（Wilson 圈）-> 非阿贝尔 -> 路径积分/正则两条leg -> 标架化。读正文时可把每一节对回这四句。
\begin{yizhu}[物理图像]
原文摘要仅四句，却已给出全文骨架：先做 $\mathrm{U}(1))$ 现象学（\S\ref{sec:2}），再做两种量子化（\S\ref{sec:3} 路径积分、\S\ref{sec:4} 正则），最后落在标架化这个技术核心上。阅读时宜把每一节对回这四句，否则容易在 \S\ref{sec:3.3} 的指标定理细节中丢失主线。
\end{yizhu}

\section*{译者说明}
\addcontentsline{toc}{section}{译者说明}

\begin{enumerate}[leftmargin=1.6em,itemsep=4pt,topsep=4pt]
  \item \textbf{底本。}David Grabovsky，\emph{Chern--Simons Theory in a Knotshell}，2022 年 9 月 14 日，38 页（pdf\TeX{} 排版）。这是一篇以 Witten 1989 年关于量子场论与 Jones 多项式的论文 \cite{witten1989} 为主线的\textbf{教学性综述}，作者写作时并非 Witten 本人；正文中凡出现「Witten 表明/证明了……」处，指的都是被综述的 Witten 工作，而非本文作者。
  \item \textbf{编号。}原文带编号的公式共 74 个（(2.1)--(2.31)、(3.1)--(3.32)、(4.1)--(4.7)、(5.1)--(5.4)）。中译本逐个保留原编号（用 \verb|\tag| 锁定），正文交叉引用一律用 \verb|\eqref|，故中译本公式号与原文逐号对应，便于对照阅读。章节号由 \LaTeX{} 自动生成，亦与原文一致。
  \item \textbf{译注。}橙色方框为译者所加的学术批注，分「易错点」「含糊点」「物理图像」「原文疑误」四类；每条批注在源码中另有一行 \verb|% [Tutor Note: 类型] 摘要| 注释，便于检索。在导言区把 \verb|\showyizhutrue| 改为 \verb|\showyizhufalse| 可一键关闭全部译注；蓝色「本节重点」方框由 \verb|\showkeytrue| 控制。两开关全关后，译本即为纯译文。
  \item \textbf{插图。}原文共 5 张位图插图（p.5、p.10、p.23、p.24、p.31），均从原 PDF 无损提取（非屏幕截图），存于 \verb|figures/|，按原位置插入，图号与原文一致。
  \item \textbf{编译。}\verb|xelatex main.tex| 连跑两遍。中文由 \verb|ctex| 宏包处理，故\textbf{必须用 Xe\LaTeX}；如无 Fandol 字体，可把导言区的 \verb|fontset=fandol| 改为 \verb|fontset=windows|（Windows 自带宋体/黑体）。
  \item \textbf{存疑。}原文中经核对属笔误或论证不严处，正文一律照录，另加「原文疑误」译注说明，并在 \verb|README.md| 的疑点清单中登记；中译本不作静默修正。
\end{enumerate}

\newpage
\tableofcontents
\newpage
